\documentclass[a4paper,11pt]{article}

% --- 言語・フォント設定 (LuaLaTeX用) ---
\usepackage{luatexja}
\usepackage{luatexja-fontspec}
\usepackage[japanese,english]{babel}

% --- 数学パッケージ ---
\usepackage{amsmath, amssymb, amsthm}
\usepackage{mathtools}
\usepackage{mathrsfs}

% --- グラフィックス・図式 ---
\usepackage{tikz}
\usetikzlibrary{shapes, backgrounds, positioning, patterns, calc, arrows.meta, fit, shadows, trees}
\usepackage{graphicx}

% --- デザイン・枠 ---
\usepackage[most]{tcolorbox}
\usepackage{geometry}
\geometry{top=2.5cm, bottom=2.5cm, left=2.5cm, right=2.5cm}

% --- ハイパーリンク ---
\usepackage{hyperref}
\hypersetup{
    colorlinks=true,
    linkcolor=blue,
    filecolor=magenta,      
    urlcolor=cyan,
}

% --- 定理環境の定義 ---
\theoremstyle{definition}
\newtheorem{definition}{定義}[section]
\newtheorem{theorem}{定理}[section]
\newtheorem{lemma}{補題}[section]
\newtheorem{example}{例}[section]
\newtheorem{remark}{注釈}[section]

% --- カスタムコマンド ---
\newcommand{\LeanCode}[1]{\texttt{#1}}
\newcommand{\Filtration}{\mathbb{F}}
\newcommand{\Algebra}[1]{\mathcal{F}_{#1}}
\newcommand{\StopT}{\tau}

% --- タイトル情報 ---
\title{\textbf{離散確率論の時系列構造：フィルトレーションと停止時間}\\
\large --- Lean 4 による形式化と情報の流れ ---}

\author{
    \textbf{TeX構成・解説}: Gemini (Google) \\
    \textbf{Lean実装・原案}: CODEX (OpenAI)
}
\date{\today}

\begin{document}

\maketitle

\begin{abstract}
本稿では、`DecidableFiltration.lean` で定義された「離散フィルトレーション（Decidable Filtration）」と「停止時間（Stopping Time）」について解説する。これらは確率過程論において「時間の経過に伴う情報の蓄積」と「未来の情報をカンニングせずに決断を下す戦略」を数学的に定式化したものである。
\end{abstract}

\tableofcontents

\section{はじめに}
確率論に「時間」の概念を導入するには、単一の代数ではなく、時間の経過とともに変化する代数の族を考える必要がある。時間が進むにつれて、私たちはより多くの情報を手に入れ、より細かい事象を区別できるようになる。

\section{離散フィルトレーション (Decidable Filtration)}

\subsection{定義}
フィルトレーションとは、単調増加する代数の列である。

\begin{tcolorbox}[colback=blue!5!white, colframe=blue!75!black, title=定義 1: 離散フィルトレーション]
有限標本空間 $\Omega$ 上のフィルトレーション $\Filtration = (\mathcal{F}_t)_{t=0}^N$ とは、以下の条件を満たす有限代数の列である：
\begin{enumerate}
    \item \textbf{時間枠 (Time Horizon)}: 有限の最終時刻 $N$ を持つ。
    \item \textbf{単調性 (Monotonicity)}: 任意の時刻 $s \le t$ に対して、
    \[ \mathcal{F}_s \subseteq_a \mathcal{F}_t \]
    すなわち、過去に観測可能だった事象は、現在も観測可能である（情報は失われない）。
\end{enumerate}
\end{tcolorbox}

\subsection{視覚的解釈：情報の細分化}
代数の包含関係 $\mathcal{F}_s \subseteq \mathcal{F}_t$ は、空間の分割が「細分化（Refinement）」されていくことに対応する。

\begin{figure}[h]
    \centering
    \begin{tikzpicture}
        % Time t=0
        \begin{scope}[xshift=0cm]
            \draw[thick, fill=gray!10] (-2,-2) rectangle (2,2);
            \node at (0, 2.3) {$t=0$ ($\mathcal{F}_0$)};
            \node at (0,0) {Unknown};
        \end{scope}
        
        % Arrow
        \draw[->, thick] (2.5, 0) -- (3.5, 0);

        % Time t=1
        \begin{scope}[xshift=6cm]
            \draw[thick, fill=gray!10] (-2,-2) rectangle (2,2);
            \draw[thick] (0,-2) -- (0,2); % Split
            \node at (0, 2.3) {$t=1$ ($\mathcal{F}_1$)};
            \node at (-1,0) {Case A};
            \node at (1,0) {Case B};
        \end{scope}

        % Arrow
        \draw[->, thick] (8.5, 0) -- (9.5, 0);

        % Time t=2
        \begin{scope}[xshift=12cm]
            \draw[thick, fill=gray!10] (-2,-2) rectangle (2,2);
            \draw[thick] (0,-2) -- (0,2); % Previous Split
            \draw[thick] (-2,0) -- (0,0); % New Split left
            \draw[thick] (0,1) -- (2,1);  % New Split right
            \node at (0, 2.3) {$t=2$ ($\mathcal{F}_2$)};
            \node at (-1,1) {A1};
            \node at (-1,-1) {A2};
            \node at (1,1.5) {B1};
            \node at (1,-0.5) {B2};
        \end{scope}
    \end{tikzpicture}
    \caption{フィルトレーションのイメージ。時間が経つにつれて、世界（$\Omega$）の区別が細かくなっていく。}
\end{figure}

\section{停止時間 (Stopping Time)}

\subsection{定義}
停止時間とは、いつプロセスを停止するか（あるいは行動を起こすか）を決める確率変数 $\tau: \Omega \to \mathbb{N}$ である。ただし、「未来の情報をカンニングしてはならない」という制約がある。

\begin{tcolorbox}[colback=green!5!white, colframe=green!75!black, title=定義 2: 停止時間]
確率変数 $\tau: \Omega \to \mathbb{N}$ がフィルトレーション $\Filtration$ に関する停止時間であるとは、任意の時刻 $t$ において、
\[ \{ \omega \in \Omega \mid \tau(\omega) \le t \} \in \mathcal{F}_t \]
が成り立つことをいう。
\end{tcolorbox}

これは、「時刻 $t$ 時点で、『もう停止したかどうか』はその時点の情報 $\mathcal{F}_t$ だけで判定できなければならない」ことを意味する。

\subsection{視覚的解釈：決定木}
停止時間は、シナリオ（$\omega$）ごとの停止するタイミングを表す。

\begin{figure}[h]
    \centering
    \begin{tikzpicture}[level distance=1.5cm,
      level 1/.style={sibling distance=3cm},
      level 2/.style={sibling distance=1.5cm}]
      
      \node[circle, draw] (root) {Start}
        child {node[circle, draw] (n1) {H}
            child {node[circle, draw] (n11) {HH}}
            child {node[circle, draw, double, fill=green!20] (n12) {HT} node[below=0.2cm, text=red] {$\tau=2$}}
        }
        child {node[circle, draw, double, fill=green!20] (n2) {T} node[below=0.2cm, text=red] {$\tau=1$}
            child {node[circle, draw, dashed] (n21) {TH}}
            child {node[circle, draw, dashed] (n22) {TT}}
        };
        
      \node[right=4cm of root, align=left, draw, drop shadow, fill=white] {
        \textbf{Example: Stop on first Tail}\\
        - If T appears at $t=1$, stop ($\tau=1$).\\
        - If H appears, continue.\\
        - If HT appears at $t=2$, stop ($\tau=2$).\\
        \\
        Decision is \textit{adapted} to the tree structure.
      };
    \end{tikzpicture}
    \caption{コイン投げにおける停止時間の例。「初めて裏が出たら停止」という戦略。破線部分は停止後の仮想的な世界線。}
\end{figure}

\section{停止時間の性質}

`DecidableFiltration.lean` では、停止時間に関する以下の代数的性質が証明されている。

\subsection{順序と有界性}
停止時間 $\tau_1, \tau_2$ に対して、点ごとの順序 $\tau_1 \le \tau_2$ を定義できる。
また、すべての $\omega$ に対して $\tau(\omega) \le N$ であるとき、$\tau$ は有界であるという。
定数停止時間（常に時刻 $c$ で止まる）は有界な停止時間の最も単純な例である。

\subsection{最小値・最大値}
2つの停止時間 $\tau_1, \tau_2$ が与えられたとき、それらの最小値と最大値もまた停止時間となる。
\begin{theorem}[停止時間の閉性]
$\tau_1, \tau_2$ が停止時間ならば、以下も停止時間である：
\begin{itemize}
    \item $\min(\tau_1, \tau_2)$：「早い方のタイミングで止まる」
    \item $\max(\tau_1, \tau_2)$：「遅い方のタイミングで止まる」
\end{itemize}
\end{theorem}

\textbf{証明の直感（minの場合）}:
$\{\min(\tau_1, \tau_2) \le t\} = \{\tau_1 \le t\} \cup \{\tau_2 \le t\}$ であり、右辺の事象はどちらも $\mathcal{F}_t$ に含まれるため、その和集合も $\mathcal{F}_t$ に含まれる（代数の性質）。

\section{実装された具体例}

\subsection{定数フィルトレーション}
情報が全く増えない状況。$\mathcal{F}_t = \mathcal{F}_0$（一定）。
この場合、停止時間は「最初から決まっている定数」のようなものしか許容されにくくなる（より正確には、$\mathcal{F}_0$-可測な定数関数）。

\subsection{2ステップ・フィルトレーション}
時刻 0 と 1 だけを考える最小の非自明な例。
\begin{itemize}
    \item $t=0$: $\mathcal{F}_0$
    \item $t=1$: $\mathcal{F}_1$ （ただし $\mathcal{F}_0 \subseteq \mathcal{F}_1$）
\end{itemize}
これは条件付き確率や、1期間モデル（ファイナンス等）の基礎となる。

\section{まとめと展望}
本稿では、情報の流れを表すフィルトレーションと、その情報に基づいて行動を決定する停止時間を定義した。
これらは、次のステップである\textbf{マルチンゲール（Martingale）}の理論---「公平な賭け」の数学的定式化---において不可欠な土台となる。特に、Optional Stopping Theorem（停止時刻定理）の計算可能な証明への道が開かれたことになる。

\end{document}